\begin{proposition}[The obstruction to a marked period lift]\hypertarget{tag-030H}{}\label{push:src:B.4.13}
\StatementRecord{push:src:B.4.13}{proved}{030H}{\SeriesRef{push:src:B.4.10}, \SeriesRef{push:src:A.1.2}, \SeriesRef{push:src:B.4.8}}

Fix $0<|\alpha|,|a|<1$, and put
\[
 H=(\C^2\setminus\{0\})/
       \langle(x,v)\mapsto(\alpha x,av)\rangle,
       \qquad E=\{x=0\}\simeq\C^*/a^\Z.
\]
Let $h$ be the automorphism induced by
$(x,v)\mapsto(c x,bv)$, where $b,c\in\C^*$.
Suppose $p:Y\to H$ is a finite composition of point blow-ups,
allowing infinitely near centres, and at least one point of $E$
is an image of its exceptional locus. If $[b]\in\C^*/a^\Z$ has
infinite order, there is no holomorphic automorphism of $Y$
which agrees, on the dense unmodified core, with the marked
automorphism $h$.

The same obstruction holds for a finite collection of such marked
modifications if the period permutes the collection and
acts as $h$ on each marked core.
More generally, if a compact smooth surface contains a preserved
marked elliptic curve $E$ of nonzero normal degree $d$, every
holomorphic automorphism restricting to translation by $b$ on $E$
must satisfy $d\,b=0$ in $E$.
The latter assertion also excludes a finite cycle of such preserved
marked cap axes when the cycle composite is an infinite-order
translation. These assertions apply to the extension
of the second cusp period to the finite Hopf
modifications in \ref{push:src:B.4.10}, for generic nonzero parameters, after every
finite base change. The period action in these assertions is a holomorphic action on
the full finite modifications with the specified marked core.
\end{proposition}


\begin{proof}
\begin{enumerate}[label=\arabic*.,leftmargin=*]
\item Let $S\subset H$ be the image of the exceptional curves of $p$.
It is finite. Repeated infinitely near blow-ups change the curves
over a point but do not make $S$ infinite.
If $\widetilde h$ were a lift, the maps
\[
 p\widetilde h,\ hp:Y\longrightarrow H
\]
would agree on the dense core and hence everywhere, by the identity
principle for holomorphic maps to the separated complex surface $H$.
The image under $\widetilde h$ of a curve contracted by $p$ is
therefore again contracted by $p$. Applying the same argument to
$\widetilde h^{-1}$ proves that the total exceptional locus is
preserved. Thus $h(S)=S$.

\item The automorphism $h$ preserves $E$, where it is translation by
$[b]$. Consequently the nonempty finite set $S\cap E$ would be
invariant under that translation. The orbit of any one of its
points is finite, so for some $k>0$,
\[
 [b^k v]=[v],\qquad\text{equivalently } b^k\in a^\Z.
\]
This says exactly that $[b]$ is torsion, a contradiction.

\item For finitely many modifications $p_i:Y_i\to H$, a deck map has biholomorphisms
$Y_i\to Y_{\sigma(i)}$ satisfying the same core comparison.
The identity principle again gives equivariance everywhere.
The union of the finitely many exceptional image sets is
$h$-invariant; its intersection with $E$ gives the same contradiction.
Equivalently, around a cycle of length $\ell$ the map $h^\ell$
would lift to a single modification, while $[b^\ell]$ still has
infinite order.

\item For the more general assertion, an automorphism preserving $E$ induces through its derivative in the normal
direction an isomorphism
\[
 t_b^*N_{E/S}\simeq N_{E/S}.
\]
For a line bundle $N$ of degree $d$ on an elliptic curve, the
homomorphism $b\mapsto t_b^*N\otimes N^{-1}$ is multiplication
by $-d$ under the usual identification $\operatorname{Pic}^0(E)=E$.
Indeed write $N=\OO_E(\sum_j n_j[p_j])$ with
$\sum_j n_j=d$; the Abel--Jacobi class of its translated divisor
minus the original one is $-d b$.
Thus preservation of $N$ forces $d b=0$.
Taking a power proves the finite-cycle version, since its composite
preserves one axis and its nonzero-degree normal line.
In particular a nonzero-degree parabolic axis cannot carry the
marked generic period translation even if no finite Hopf
modification is specified. The original finite-support argument
also covers modifications where that normal-degree test alone
would give no information.

\item The infinite-order translation at the cusp is computed as follows.
The periods of \ref{push:src:A.1.2} have the form
\[
 g_1:(Z,W)\longmapsto(q c_1 Z,c_2W),\qquad
 g_2:(Z,W)\longmapsto(c_3Z,r^{12}q c_4 W),
\]
where the $c_i$ are nonvanishing holomorphic units near the cusp.
Choose $W$ or $W^{-1}$ so that
$a=c_2^{\varepsilon}$ has modulus less than one,
$\varepsilon\in\{1,-1\}$.
This is possible at the marked cusp because its limiting
unit has modulus different from one, as computed in \ref{push:src:B.4.8}.
For $q\ne0$ small, the partial $g_1$ quotient of $(\C^*)^2$
is the dense core of the diagonal Hopf surface with multipliers
$(q c_1,a)$. The second period induces the above $h$, with
\[
 b(q,r)=(r^{12}q c_4(q,r))^{\varepsilon}.
\]
Fix a sufficiently small nonzero $q$. For each $k>0$ and
$j\in\Z$, the equation $b(q,r)^k=a(q)^j$ has a discrete
solution set in the punctured $r$-disc: its left-hand side has
a zero or pole of nonzero order at $r=0$, while the right-hand
side is a nonzero constant in $r$.
Hence outside a countable union of discrete sets, $[b(q,r)]$
has infinite order. Under a finite base change in $r$ the
nonzero order is multiplied by its ramification index and
the same conclusion holds. A finite base change in $q$ does
not affect this fixed-nonzero-$q$ argument.

\item To compare with \ref{push:src:B.4.10}, write $q c_1=\kappa(s)s^n$ after
the chosen finite cover in the cusp variable. Its generic Kato
fibre is precisely a finite Hopf modification with nonempty
axis centre $[-s]$. A holomorphic extension of the second
period to that whole fibre would therefore contradict the result
just proved. Allowing finitely many translated copies does not
remove the contradiction, by the finite-collection assertion.
\end{enumerate}
\end{proof}